% !TEX root = ../main.tex
\chapter*{Introducción}
\addcontentsline{toc}{chapter}{Introducción}
\markboth{Introducción}{Introducción}

La patología digital ha transformado la adquisición, conservación y análisis de
las láminas histológicas al sustituir la observación exclusivamente microscópica
por imágenes digitales de alta resolución que pueden almacenarse, consultarse y
compartirse~\cite{pallua_future_2020,song_artificial_2023}. Este cambio hace
posible aplicar métodos computacionales a problemas diagnósticos en los que la
morfología del tejido constituye una fuente principal de información. El cáncer,
entendido como un conjunto de enfermedades caracterizadas por el crecimiento y
la diseminación descontrolados de células anormales~\cite{nci_what_is_cancer}, es
un ámbito de especial interés. En Colombia fue la segunda causa de muerte durante
el periodo 2007--2013, con aproximadamente 33.538 defunciones anuales, equivalentes
al 17,1\% de la mortalidad del país~\cite{pardo_c_atlas_2017}.

En muchos procesos diagnósticos se extrae una biopsia, se prepara el tejido sobre
un portaobjetos y un patólogo examina sus características microscópicas; sin
embargo, la biopsia no es el único procedimiento diagnóstico y su necesidad
depende del caso clínico~\cite{nci_cancer_diagnosis}. La interpretación
histopatológica tampoco es completamente uniforme. En carcinoma de ovario, la
comparación entre el grado consignado en registros poblacionales y una revisión
central de 664 tumores produjo un coeficiente kappa de 0,25
~\cite{matsuno_agreement_2013}. En cuatro ensayos del \emph{Cooperative Breast
Cancer Tissue Resource}, con tamaños de muestra de 10, 10, 19 y 10 casos, la
reproducibilidad del grado fue moderada, con valores kappa entre 0,50 y
0,59~\cite{meyer_breast_2005}. En próstata, la reevaluación de biopsias de 350
pacientes coincidió completamente con el informe inicial en el 76,9\% de los
puntajes de Gleason y modificó decisiones de evaluación o tratamiento en una
fracción de los pacientes~\cite{berg_prostate_2011}. Estos resultados ilustran
que la variabilidad de interpretación debe considerarse al diseñar y evaluar
herramientas de apoyo.

El análisis computacional de una WSI presenta, además, dificultades de escala.
Una sola lámina puede aproximarse a \(100\,000\times100\,000\) píxeles, por lo
que no suele procesarse como una imagen convencional en memoria
~\cite{wang_managing_2012}. A la vez, la evidencia relevante puede ocupar solo
algunas regiones del tejido. Por ello, los sistemas de patología computacional
dividen las láminas en parches, aprenden representaciones locales y agregan la
información necesaria para producir una predicción global
~\cite{niazi_digital_2019,song_artificial_2023}.

El entrenamiento supervisado de estos sistemas requiere anotaciones producidas
por especialistas. Su costo depende de la tarea y del nivel de detalle; en un
proyecto de anotación exhaustiva de 200 WSI de colon y piel se registraron
tiempos desde 45 segundos hasta más de 360 minutos por lámina
~\cite{lindman_annotations_2019}. La disponibilidad simultánea de numerosas
imágenes sin anotación y un subconjunto menor con etiquetas motiva entonces un
esquema semisupervisado: primero se aprende una representación sin utilizar las
etiquetas diagnósticas y después se evalúa esa representación mediante tareas
supervisadas con anotaciones limitadas.

Para estudiar si una restricción geométrica puede mejorar ese aprendizaje, este
trabajo compara un autoencoder variacional convolucional base con otro
equivariante a la acción continua del grupo de rotaciones planas
\(\mathrm{SO}(2)\). La elección se apoya en que la orientación global de las
estructuras histopatológicas no determina por sí sola su clase y en que las
arquitecturas equivariantes pueden explotar esta regularidad
~\cite{veeling_rotation_2018}. Los beneficios teóricos de imponer una simetría
dependen, no obstante, de que la distribución y la tarea satisfagan los supuestos
correspondientes~\cite{elesedy_provably_2021}; por esta razón la tesis no supone
de antemano una superioridad empírica. El alcance comprende rotaciones continuas,
pero no reflexiones ni una extensión adicional de la equivarianza traslacional
propia de la convolución.

En resumen, se implementaron y evaluaron dos modelos de representación bajo un
protocolo pareado. La comparación considera reconstrucción, clasificación a nivel
de WSI, segmentación de WSI a resolución de parche y respuesta del espacio latente
a rotaciones continuas. La segmentación no se interpreta como una cobertura
exhaustiva de la lámina porque las máscaras disponibles contienen regiones sin
anotar que no pueden utilizarse como etiquetas negativas.
